% !TEX root = ../main.tex
\clearpage
\section{Opciones metodológicas para la evaluación}
\label{app:evaluacion-piloto}
\textbf{Propuesta para discusión.} Este apéndice concreta alternativas para la secuencia de la sección~\ref{sec:piloto}. La elección del diseño dependerá del objetivo institucional, los servicios disponibles y los recursos; no constituye un protocolo aprobado.

\subsection{Calidad de las notas sobre casos comunes}
Se propone evaluar los productos con las mismas consultas y el mismo contexto clínico, separando los casos de ajuste de plantillas de los reservados para evaluación. Las diferencias entre carga de archivos, reproducción de audio y captura en directo deben documentarse. Los trabajos de Biro y Ha ofrecen procedimientos con consultas estandarizadas y revisión clínica que pueden adaptarse; sus muestras exploratorias no determinan el tamaño necesario para este hospital~\cite{NNWFT2JX,PRTISMYF}.

La referencia combinaría audio, contexto disponible y valoración independiente. La nota original aporta información, pero también puede contener omisiones. Se propone contar con dos revisores, ocultar la identidad del proveedor cuando sea viable y resolver los desacuerdos. La valoración distinguiría errores clínicos, omisiones, gravedad y preferencias de estilo; la similitud automática entre textos no sustituye el juicio clínico~\cite{NNWFT2JX,4IHZ7QFQ}.

\subsection{Diseño para evaluar el efecto asistencial}
Un banco de consultas permite comparar borradores, pero no estimar por sí solo el ahorro de trabajo real. Para este segundo objetivo se plantean tres alternativas:
\begin{small}
\begin{longtable}{L{.25\linewidth}L{.31\linewidth}L{.34\linewidth}}
\toprule
Diseño & Utilidad & Limitación principal \\
\midrule\endhead
Antes/después con los mismos profesionales & Explorar factibilidad y cambios iniciales. & Confunde el efecto del producto con aprendizaje y cambios de actividad. \\
Grupo contemporáneo con documentación habitual; asignación aleatoria si es viable & Estimar el efecto frente al flujo habitual. & Requiere muestra suficiente y equilibrio entre servicios y carga inicial. \\
Incorporación escalonada o periodos alternos & Acompañar un despliegue progresivo. & Requiere controlar calendario, aprendizaje y posibles efectos persistentes entre periodos. \\
\bottomrule
\end{longtable}
\end{small}
Estas alternativas son propuestas de diseño, no una clasificación validada de protocolos para escribas. El ensayo de UCLA aporta un ejemplo de comparación aleatorizada con dos productos y control~\cite{PR93N8FN}. Para evitar confundir producto y especialidad, convendría comparar dentro de servicios o equilibrar la asignación.

\clearpage
\subsection{Resultados y decisiones pendientes}
Como objetivo principal inicial se plantea el \textbf{tiempo total hasta completar la documentación}, condicionado a una calidad clínica aceptable. Si la prioridad fuera bienestar profesional, deberían ajustarse el resultado principal y el seguimiento. Las medidas complementarias incluirían:
\begin{itemize}
\item \textbf{Calidad:} errores relevantes, omisiones y errores residuales tras revisión, con denominadores y gravedad definidos.
\item \textbf{Trabajo real:} revisión, edición dentro y fuera del escriba, incorporación a HCIS y trabajo fuera de horario. Los registros de una aplicación pueden omitir correcciones realizadas en otra~\cite{C7S8N9NA}.
\item \textbf{Adopción y experiencia:} uso sobre consultas elegibles, rechazos, abandono, carga percibida y experiencia del paciente. Las expectativas previas y la experiencia durante el uso deben analizarse por separado~\cite{PR93N8FN,54NFKG36}.
\item \textbf{Operación y coste:} fallos, latencia, recuperación y coste por consulta con documentación completada.
\end{itemize}

Antes del inicio deberán definirse condiciones de tratamiento del conjunto de prueba, formación, versiones, criterios de suspensión y cambios que exijan repetir la evaluación. La muestra responderá al objetivo y a la precisión deseada, considerando la agrupación de consultas por profesional. La guía \fuente{https://discovery.ucl.ac.uk/10152188/3/Supplementary_appendix_1_clean.pdf}{DECIDE-AI} aporta criterios de reporte sobre evaluación inicial y factores humanos; su adaptación al escriba documental sería metodológica, pues su ámbito original es el apoyo a decisiones clínicas. La autorización para preparar datos de evaluación no equivale a autorizar su transferencia a proveedores~\cite{rgpd2016}.
